\begin{resumo}[Resumo]
\begin{otherlanguage*}{brazil}
Bolhas de Plasma Equatoriais (EPBs) desencadeiam cintilações ionosféricas severas, frequentemente interrompendo a comunicação e a navegação por satélite. Durante anos, os pesquisadores confiaram na detecção manual em imagens de luminescência atmosférica de todo o céu. Embora fundamental, essa abordagem exige um trabalho imenso e permanece sensível ao ruído ambiental, apontando para uma clara necessidade de automação. O aprendizado profundo certamente abriu novas portas aqui, mas o uso de imagens ópticas brutas diretamente em redes complexas tende a aumentar os custos computacionais e amplificar artefatos de fundo. Em resposta, desenvolvemos uma arquitetura de aprendizado de máquina baseada na física, adaptada especificamente para a detecção de EPB. O processo começa com um pré processamento específico do domínio. Ao modelar e subtrair a iluminação de fundo difusa, juntamente com a Equalização de Histograma Adaptativa Limitada por Contraste (CLAHE), nós acentuamos naturalmente as estruturas físicas do plasma. Em vez de achatar essas imagens e perder seu contexto espacial, aplicamos a Análise de Componentes Principais Bidimensional (2DPCA) para extrair características diretamente. Essa etapa reduz substancialmente a dimensionalidade, mantendo intactas as correlações espaciais bidimensionais. Um classificador ensemble stacking, combinando os algoritmos Support Vector, XGBoost e Random Forest, avalia essas características. Também incorporamos o SMOTE para gerenciar o desequilíbrio de classe inerente entre céus limpos e eventos de bolhas. Ao longo de uma validação cruzada de 5 dobras, o modelo atingiu uma Acurácia média e um F1 Score de 99.80\%. Além de apenas métricas, as explicações aditivas de SHapley (SHAP) provaram que o algoritmo realmente foca na morfologia física das depleções, ignorando ruídos ambientais aleatórios. Quando testado em sequências noturnas contínuas, o sistema manteve forte estabilidade virtualmente sem falsos positivos, oferecendo uma ferramenta rápida e altamente interpretável para o monitoramento do clima espacial em tempo real.

\vspace{\onelineskip}
\noindent Palavras chave: Bolhas de Plasma Equatoriais. EPB. Aprendizado de Máquina. Clima Espacial. Imageador de Todo o Céu. 2DPCA. Ensemble Stacking. Irregularidades Ionosféricas.
\end{otherlanguage*}
\end{resumo}
